\documentclass[conference]{IEEEtran}
\IEEEoverridecommandlockouts

\usepackage{cite}
\usepackage{amsmath,amssymb,amsfonts}
\usepackage{algorithmic}
\usepackage{graphicx}
\usepackage{textcomp}
\usepackage{xcolor}
\usepackage{booktabs}
\usepackage{url}
\def\BibTeX{{\rm B\kern-.05em{\sc i\kern-.025em b}\kern-.08em
    T\kern-.1667em\lower.7ex\hbox{E}\kern-.125emX}}

\title{Machine Learning for Modelling and Control of an Unbalanced Disk\\
Group 17 Design Project}

\author{\IEEEauthorblockN{TODO Student 1 Name, TODO Student 2 Name, TODO Student 3 Name}
\IEEEauthorblockA{\textit{5SC28 Machine Learning for Systems and Control}\\
\textit{Eindhoven University of Technology}\\
TODO student numbers}
}

\begin{document}

\maketitle

\begin{abstract}
% Target length: 120--180 words.
% Write this last.
This report presents the Group 17 design project on machine learning for modelling and control of an unbalanced disk.
The modelling part identifies the system dynamics from measured data, where the input is the applied motor voltage and the output is the measured disk angle.
Both Gaussian Process and Artificial Neural Network models are developed and evaluated on prediction and simulation tasks.
The control part studies reinforcement learning policies for swing-up and reference tracking.
The methods are compared using validation errors, hidden-test output files, simulation performance, and, if available, real setup experiments.
\end{abstract}

\begin{IEEEkeywords}
unbalanced disk, system identification, Gaussian Process, Artificial Neural Network, LSTM, Q-learning, model internalization, reference tracking
\end{IEEEkeywords}

% =========================================================
% WRITING PLAN
% =========================================================
% Required report format from announcement:
% - PDF.
% - IEEE 2-column conference format using IEEEtran.
% - 6--12 pages.
% - Include all names and student numbers in the author block.
% - Remove all writing-plan comments before final PDF export.
% - IEEE style reminders: define abbreviations the first time they appear,
%   cite figures after they are mentioned, place figure captions below figures,
%   place table captions above tables, and use SI units.
%
% Group zip must include:
% - 1 report PDF.
% - 1 folder with requested prediction/simulation files.
% - 1 folder with relevant code and additional material.
%
% Requested test files:
% - at least one ANN prediction result.
% - at least one ANN simulation result.
% - at least one GP prediction result.
% - at least one GP simulation result.
%
% Max-grade strategy:
% - Do not only show final plots. Explain choices, tuning, failures, and limitations.
% - For every method, include: theory link, implementation choice, validation method, final result, critical discussion.
% - Make prediction vs simulation clear. The lecturers care about this.
% - Mention that angular velocity was not used in the modelling part.
% - Leave real setup results honest: if not tested, say so; if tested, compare with simulation.

\section{Introduction}
% Page budget: about 0.5 page.
% Max-grade purpose: clearly introduce the problem and show that the report follows the assignment objectives.

\subsection{Motivation and Problem}
The unbalanced disk is a nonlinear dynamical system.
The project goal is to use machine learning methods from the course to model the system dynamics and to learn controllers for swing-up and reference tracking.

Write:
\begin{itemize}
    \item what the unbalanced disk is;
    \item why it is a useful control benchmark;
    \item why nonlinear modelling and reinforcement learning are relevant;
    \item what ``successful'' modelling and control mean in this project.
\end{itemize}

\subsection{Assignment Objectives}
The assignment is split into three main objectives:
\begin{enumerate}
    \item identify the system dynamics using GP and ANN models;
    \item learn a swing-up policy;
    \item extend the policy to reference tracking around the upright position.
\end{enumerate}

State the exact modelling variables:
\begin{equation}
u(k) = \text{applied motor voltage}, \qquad
\theta(k) = \text{measured disk angle}.
\end{equation}

Important sentence to include:
For the modelling task, angular velocity was not used, following the assignment requirement.

\subsection{Group 17 Contributions}
% Keep this short in the group report.
% The individual reflection is submitted separately and has max 500 words.

\begin{itemize}
    \item ANN modelling: TODO name and one-sentence contribution.
    \item GP modelling: TODO name and one-sentence contribution.
    \item Control/simulation: TODO name and one-sentence contribution.
    \item Real setup test, if performed: TODO names and date.
\end{itemize}

\section{System, Data, and Evaluation Tasks}
% Page budget: about 0.75--1 page.
% Max-grade purpose: make the experimental setup and evaluation rules unambiguous.

\subsection{System Description}
Describe the measured disk angle, the motor voltage, the sampling of the data, and the goal of reaching the upright position.
Mention that the output is noisy.

Suggested figure:
\begin{itemize}
    \item Figure 1: photo or schematic of the unbalanced disk setup.
\end{itemize}

\subsection{Available Data}
Describe:
\begin{itemize}
    \item training/validation/test data from the provided files;
    \item hidden prediction test file;
    \item hidden simulation test file;
    \item how the group split the known data for validation before generating hidden-test outputs.
\end{itemize}

\subsection{Prediction Versus Simulation}
This part is important for theory discussion.

Prediction uses measured past outputs:
\begin{equation}
\hat{\theta}(k) = f(u(k-1), \ldots, \theta(k-1), \ldots).
\end{equation}

Simulation feeds the model's own previous predicted outputs back into the model.
This is harder because errors can accumulate over time.

State clearly:
\begin{itemize}
    \item prediction tests one-step-ahead accuracy;
    \item simulation tests long-term dynamical consistency;
    \item both are required by the assignment.
\end{itemize}

\subsection{Evaluation Metrics}
Use RMSE in radians and degrees:
\begin{equation}
\mathrm{RMSE} =
\sqrt{\frac{1}{N}\sum_{k=1}^{N}(\theta(k)-\hat{\theta}(k))^2}.
\end{equation}

Mention any extra control metrics:
\begin{itemize}
    \item success rate;
    \item mean absolute tracking error;
    \item percentage of time near upright/reference;
    \item control voltage behaviour.
\end{itemize}

\section{System Dynamics Modelling}
% Page budget: about 2--3 pages.
% Max-grade purpose: cover both GP and ANN theory, design choices, tuning, and final results.

\subsection{System Identification Approach}
Connect to Lecture 1 and the standard system-identification cycle \cite{course5sc28,ljung1999system}:
\begin{itemize}
    \item choose a model structure;
    \item estimate parameters or train the model;
    \item validate using unseen data;
    \item tune hyperparameters;
    \item compare prediction and simulation;
    \item discuss limitations and uncertainty.
\end{itemize}

Explain why nonlinear models are needed for the unbalanced disk.

\subsection{Gaussian Process Model}
% Fill this using the GP teammate's folder.
% Max-grade purpose: do not only report the score. Explain why the GP model makes sense.

\subsubsection{Model Structure}
TODO describe the GP input-output structure.
For example, if GP-NARX is used:
\begin{equation}
\hat{\theta}(k) =
f(u(k-1), \ldots, u(k-n_u), \theta(k-1), \ldots, \theta(k-n_y)).
\end{equation}

Mention:
\begin{itemize}
    \item delays used;
    \item whether the model predicts angle directly or a difference/residual;
    \item whether sparse GP is used.
\end{itemize}

\subsubsection{Kernel and Hyperparameters}
Explain the theory using the GP lecture material and a standard GP reference \cite{course5sc28,rasmussen2006gpml}:
\begin{itemize}
    \item GP defines a distribution over functions;
    \item mean function and covariance/kernel;
    \item kernel choice controls smoothness and similarity;
    \item hyperparameters are learned/tuned from data;
    \item sparse GP reduces computational cost if used.
\end{itemize}

\subsubsection{GP Results}
Add:
\begin{itemize}
    \item prediction RMSE;
    \item simulation RMSE;
    \item final hidden-test files generated;
    \item plot of validation simulation;
    \item short discussion of uncertainty/limitations.
\end{itemize}

\subsection{Artificial Neural Network Models}
% This is your ANN part. It should be detailed enough to defend.

\subsubsection{Simple ANN: NARX Feedforward Network}
The simple ANN uses delayed voltages and delayed measured angles.
Neural networks are a standard tool for nonlinear dynamical-system identification \cite{narendra1990identification}.
\begin{equation}
\hat{\theta}(k) =
f_{\mathrm{ANN}}(u(k-1), \ldots, u(k-n_u),
\theta(k-1), \ldots, \theta(k-n_y)).
\end{equation}

Write:
\begin{itemize}
    \item this is the simple model required by section 4.1;
    \item it follows the NARX idea from the practical sessions;
    \item it uses only motor voltage and disk angle;
    \item it is easy to train but has limited memory because delays are fixed manually.
\end{itemize}

Theory to mention:
\begin{itemize}
    \item feedforward neural networks approximate nonlinear functions;
    \item training uses backpropagation;
    \item overfitting is controlled by validation/early stopping or selected epochs.
\end{itemize}

\subsubsection{Advanced ANN: GRU and LSTM Recurrent Networks}
The advanced ANN models use a short sequence.
GRU and LSTM networks are recurrent neural networks designed to learn sequence dependencies.
GRU is a simpler gated recurrent model \cite{cho2014gru}, while LSTM was designed to reduce the vanishing-gradient problem \cite{hochreiter1997lstm}.
\begin{equation}
[(u(k-H),\theta(k-H)), \ldots, (u(k-1),\theta(k-1))]
\rightarrow \hat{\theta}(k).
\end{equation}

Write:
\begin{itemize}
    \item GRU and LSTM are the advanced ANN architectures used for section 4.1;
    \item recurrent networks are suitable for time series;
    \item gated recurrent networks help with sequence dependencies and vanishing gradients;
    \item it still does not use angular velocity.
\end{itemize}

\subsubsection{Training, Validation, and Tuning}
Include exact settings:
\begin{itemize}
    \item training/validation/test split;
    \item normalization method;
    \item optimizer and learning rate;
    \item input delay/output delay for NARX;
    \item history length and hidden size for GRU/LSTM;
    \item number of epochs;
    \item tuning settings tried.
\end{itemize}

Suggested table:
\begin{table}[h]
\centering
\caption{ANN hyperparameter settings.}
\begin{tabular}{lccc}
\toprule
Model & Memory & Hidden size & Epochs \\
\midrule
NARX ANN & TODO & TODO & TODO \\
GRU ANN & TODO & TODO & TODO \\
LSTM ANN & TODO & TODO & TODO \\
\bottomrule
\end{tabular}
\end{table}

\subsection{Modelling Results and Comparison}
% Max-grade purpose: compare all models and explain why the final selected files were chosen.

\begin{table}[h]
\centering
\caption{Validation results for system dynamics models.}
\begin{tabular}{lcc}
\toprule
Model & Prediction RMSE & Simulation RMSE \\
\midrule
Simple NARX ANN & TODO & TODO \\
Advanced GRU ANN & TODO & TODO \\
Advanced LSTM ANN & TODO & TODO \\
Gaussian Process & TODO & TODO \\
\bottomrule
\end{tabular}
\end{table}

Figures to include:
\begin{itemize}
    \item prediction comparison plot;
    \item simulation comparison plot;
    \item ANN tuning plot;
    \item GP simulation plot;
    \item error bar plot comparing ANN and GP.
\end{itemize}

Discussion points:
\begin{itemize}
    \item which model is best for prediction;
    \item which model is best for simulation;
    \item why simulation errors are usually larger;
    \item whether the advanced ANN improved over the simple ANN;
    \item whether the GP has advantages in uncertainty or data efficiency;
    \item final files selected for submission.
\end{itemize}

\section{Policy Learning for Swing-Up}
% Page budget: about 1.5--2 pages.
% Max-grade purpose: show RL theory, reward design, training design, and actual performance.

\subsection{Control Objective}
The policy should move the disk from the stable bottom position to the upright position and keep it there.
Use reinforcement-learning theory from the course and Sutton and Barto \cite{course5sc28,sutton2018reinforcement}.
Mention constraints such as voltage limits and the nonlinear nature of the task.

\subsection{Q-Learning / DQN Method}
TODO fill using the control folder.

Theory to include:
\begin{itemize}
    \item Markov decision process: state, action, reward, transition;
    \item value function or Q-function;
    \item exploration versus exploitation;
    \item why DQN is useful for high-dimensional state information, if used \cite{mnih2015human}.
\end{itemize}

Implementation details:
\begin{itemize}
    \item state variables used by the controller;
    \item action space / voltage actions;
    \item reward function;
    \item training episodes;
    \item neural network architecture, if DQN;
    \item final selected controller.
\end{itemize}

\subsection{Actor-Critic or Model Internalization Method}
TODO fill using the hybrid/model-internalization folder.

Theory to include:
\begin{itemize}
    \item difference between value-based and policy-based methods \cite{sutton1999policy};
    \item actor and critic roles, or model internalization concept;
    \item why this method can improve over basic Q-learning;
    \item robustness testing.
\end{itemize}

\subsection{Swing-Up Simulation Results}
Suggested table:
\begin{table}[h]
\centering
\caption{Swing-up controller simulation results.}
\begin{tabular}{lccc}
\toprule
Controller & Success rate & Upright error & Notes \\
\midrule
Q-learning/DQN & TODO & TODO & TODO \\
Model internalization & TODO & TODO & TODO \\
\bottomrule
\end{tabular}
\end{table}

Figures to include:
\begin{itemize}
    \item disk angle over time;
    \item control voltage over time;
    \item training return/loss curve;
    \item error near upright.
\end{itemize}

\section{Reference Tracking}
% Page budget: about 1 page.
% Max-grade purpose: show the final extension beyond swing-up.

\subsection{Objective}
The final controller should use one policy to swing up and then track a reference around the upright position.

\subsection{Method}
TODO describe the reference tracking policy:
\begin{itemize}
    \item reference signal;
    \item state representation;
    \item reward function;
    \item training setup;
    \item how swing-up and tracking are combined.
\end{itemize}

\subsection{Simulation Results}
Suggested table:
\begin{table}[h]
\centering
\caption{Reference tracking simulation results.}
\begin{tabular}{lccc}
\toprule
Controller & Success rate & Tracking RMSE & Notes \\
\midrule
Reference tracking policy & TODO & TODO & TODO \\
\bottomrule
\end{tabular}
\end{table}

Figures:
\begin{itemize}
    \item measured/simulated angle and reference on same plot;
    \item tracking error;
    \item voltage input.
\end{itemize}

\section{Real Setup Experiments}
% Page budget: 0.5--1 page depending on tomorrow's results.
% Leave this section even if the real test fails. A critical discussion helps the grade.

\subsection{Test Plan}
Before testing, record:
\begin{itemize}
    \item date and setup used;
    \item controller version;
    \item number of trials;
    \item initial condition;
    \item safety constraints;
    \item success definition;
    \item logged signals.
\end{itemize}

\subsection{Experimental Results}
TODO fill after tomorrow's test.

\begin{table}[h]
\centering
\caption{Real setup experiments.}
\begin{tabular}{lccc}
\toprule
Trial & Controller & Success & Comment \\
\midrule
1 & TODO & TODO & TODO \\
2 & TODO & TODO & TODO \\
3 & TODO & TODO & TODO \\
\bottomrule
\end{tabular}
\end{table}

Discuss:
\begin{itemize}
    \item whether the real setup matched simulation;
    \item any unexpected behaviour;
    \item noise, friction, delay, saturation, or safety issues;
    \item what would be changed with more lab time.
\end{itemize}

\section{Discussion}
% Page budget: about 1 page.
% Max-grade purpose: this is where critical thinking and theory discussion show up.

\subsection{System Identification Discussion}
Compare GP and ANN using:
\begin{itemize}
    \item accuracy;
    \item simulation stability;
    \item interpretability;
    \item hyperparameter sensitivity;
    \item computation time;
    \item suitability for control.
\end{itemize}

\subsection{Control Discussion}
Compare Q-learning, model internalization, and reference tracking:
\begin{itemize}
    \item reward design sensitivity;
    \item sample efficiency;
    \item robustness;
    \item ease of tuning;
    \item simulation-to-real transfer.
\end{itemize}

\subsection{Limitations}
Be honest and specific:
\begin{itemize}
    \item noisy disk angle measurements;
    \item no angular velocity used in modelling, as required, but this may limit model information;
    \item hidden test scores are not directly known;
    \item neural networks may overfit tuning data;
    \item real setup time is limited;
    \item simulation model may not fully match physical setup.
\end{itemize}

\subsection{Course Theory Coverage}
Use this paragraph to explicitly show theory knowledge from the rubric:
\begin{itemize}
    \item dynamic model structures, NARX, prediction, and simulation;
    \item GP mean, covariance/kernel, hyperparameters, and sparse approximation if used;
    \item ANN approximation, backpropagation, regularization/overfitting, recurrent networks, vanishing gradients;
    \item RL basics: MDP, policy, value/Q-function, reward, exploration;
    \item RL advanced: DQN, actor-critic or model internalization, robustness.
\end{itemize}

\section{Conclusion}
% Page budget: about 0.3--0.5 page.

Summarize:
\begin{itemize}
    \item best GP and ANN modelling results;
    \item whether the advanced ANN improved over the simple ANN;
    \item best swing-up and tracking controller results;
    \item real setup result, if available;
    \item final lesson: systematic validation and tuning matter more than only trying a complex model.
\end{itemize}

\section*{Acknowledgment and Contributions}
% Keep short. Full individual reflection is separate.
% Include names/student numbers somewhere in the report as required by announcement.

TODO short group contribution statement.

\bibliographystyle{IEEEtran}
\bibliography{references}

\end{document}
